\documentclass[10pt,a4paper]{amsart}
\usepackage[margin=19mm]{geometry}
\usepackage{amsmath,amssymb,mathtools}
\usepackage[scr=rsfs,cal=euler]{mathalfa}
\usepackage{xfrac}
\usepackage[unicode,hidelinks,bookmarks=false]{hyperref}
\newcommand{\ie}{i.e.\ }
\newcommand{\cf}{cf.\ }
\newcommand{\resp}{respectively\ }
\newcommand{\Subsection}{\subsection}
\providecommand{\nobreakdash}{\nobreak\hbox{-}}
\setlength{\emergencystretch}{2em}
\multlinegap0pt
\usepackage{amssymb}
\usepackage{enumerate}
\usepackage{xcolor}
\newtheorem{lemma}{Lemma}
\newcommand{\R}{\mathbb{R}}
\newcommand{\Tree}{\mathcal{T}}
\newcommand{\W}{\mathcal{W}}
\newcommand{\hd}{\hspace{0.2cm}}
\newcommand{\izj}{\int_{0}^{1}}
\newcommand{\m}[1]{\mbox{#1}}
\newcommand*{\norm}[1]{\left\Vert{#1}\right\Vert}
\newcommand{\ve}{\varepsilon}
\title{The mean-field limit of non-exchangeable particle systems with non-conservative dynamics and adaptive weights}
\author{Katarzyna Ryszewska}
\date{Source: arXiv:2607.21110v2 (CC BY 4.0)}
\begin{document}
\maketitle

\noindent\emph{Source and licence.} The mean-field limit of non-exchangeable particle systems with non-conservative dynamics and adaptive weights. Version arXiv:2607.21110v2 (CC BY 4.0), distributed by arXiv under the Creative Commons Attribution 4.0 International licence, http://creativecommons.org/licenses/by/4.0/, as recorded in the arXiv licence metadata for this exact version and shown on the abstract page https://arxiv.org/abs/2607.21110v2. Full licence notice: \texttt{licenses/arxiv-2607.21110-CC-BY-4.0.txt}.\par\smallskip

\noindent\textit{Editorial scope.} Counted unit: the article's lemma taustar (source0.tex lines 1177--1189). The article's complete original proof of it is delivered with it (source0.tex lines 2094--2159, one \texttt{proof} environment of 66 lines, which the article places away from the statement). No mathematics is written, repaired or re-derived: the statement and the proof are the article's own lines, in the article's own notation, and no retained line is substituted or re-flowed.  Retained uncounted as the source-local context the proof needs: lines 581--599; lines 662--671; lines 1019--1024; lines 1044--1046.  Nothing was omitted from any retained span, so the package carries no omission record.  The retained lines cite 1 works, whose entries are transcribed verbatim from the article's own bibliography source in the e-print and delivered with short editorial labels recorded in the item's reference mapping.  No retained line contains a figure environment, an image-inclusion command or drawing code, so no image is delivered and the package declares no asset dependency.  Editorial formatting only, in addition: an AMS \texttt{amsart} preamble that restates the article's own declarations and macros used by the retained lines, namely the environments \texttt{lemma} on separate counters, together with the standard packages those lines need.  The retained environments print in this document as Lemma 1 in order of appearance, whereas the article prints the same items with the numbering of its own full document.  The complete native source of the article as distributed in the arXiv e-print for this exact version is bundled unchanged as \texttt{sources/205855/source0.tex}.
\begin{lemma}\cite[Lemma 4.7]{JPS2025}\label{impoest}
The bilinear,  bounded operator
\[
\begin{array}{ccl}
\W \times C_\zeta([0,1]) & \longrightarrow & L^\infty_\xi((0,1)),\\
(w,\phi) & \longmapsto & \int_0^1\phi(\zeta)\,w(\cdot,d\zeta)
\end{array}
\]
 extends to a bounded operator from $\W \times L^\infty_\zeta((0,1))$ to $L^\infty_\xi((0,1))$. 
Furthermore, for any $w\in \W$ and $\phi\in L^\infty((0,1))$
\[
\left\|\int_0^1 \phi(\zeta)\,w(\cdot,d\zeta)\right\|_{L^1((0,1))}\leq \|w\|_{L^\infty_\zeta \mathcal{M}_\xi}\,\|\phi\|_{L^1((0,1))}, \hd 
\left\|\int_0^1 \phi(\zeta)\,w(\cdot,d\zeta)\right\|_{L^\infty((0,1))}\leq \|w\|_{L^\infty_\xi \mathcal{M}_\zeta}\,\|\phi\|_{L^\infty((0,1))}.
\]
In addition, for any uniformly bounded sequence $(\phi_n)_{n\in \mathbb{N}} \in L^{\infty}((0,1))$ such that $\phi_n\rightarrow \phi$ in $L^1((0,1))$ and a sequence   $w_n\overset{*}{\rightharpoonup} w$  in $\W$ the following convergence holds in weak-star topology on $L^{\infty}((0,1))$
  \[
\int_0^1\phi_n(\zeta)\,w_n(\cdot,d\zeta)\overset{*}{\rightharpoonup} \int_0^1 \phi(\zeta)w(\cdot,d\zeta).
\]
\end{lemma}

\begin{lemma}\label{hmj}
Let $f \in \W((L^1\cap L^{\infty})(\R^d))$. Let us extend $f$ periodically in $\xi$, convolve with standard smoothing periodic kernel $\rho^\ve$ in $\xi$ variable and  denote $f^{\ve}(x,\xi,\eta)=\izj \rho^\ve(\xi-\xi')f(x,d\xi',\eta)$. Then $f^\ve \in L^{\infty}((0,1)^2\times \R^d)$,  $f^\ve \rightarrow f$ in $L^{1}_\xi(H^{-1}_\eta;L^{1}(\R^d))$ and in  $L^{1}_\eta(H^{-1}_\xi;L^{1}(\R^d))$.
Moreover, for $Y=L^{\infty}(\R^d)$ or $Y=\mathcal{M}(\R^d)$  there hold the estimates $\norm{f^\ve}_{\W(Y)} \leq \norm{f}_{\W(Y)}$. Finally, for any $g, g^{\ve} \in L^{1}((0,1)\times \R^{dn}) \cap L^{\infty}((0,1);L^{1}(\R^{dn})) $ such that $g^\ve \rightarrow g$ in $L^{1}((0,1)\times \R^{dn})$, we have
\[
\izj g^{\ve}(y_1,\dots,y_n,\eta) f^{\ve}(x,\cdot,d\eta) \rightarrow \izj g(y_1,\dots,y_n,\eta) f(x,\cdot,d\eta) \m{ in } L^{1}_\xi((0,1)\times \R^{d(n+1)}), 
\]
\[
 \izj g^{\ve}(y_1,\dots,y_n,\xi) f^{\ve}(x,d\xi,\cdot) \rightarrow \izj g(y_1,\dots,y_n,\xi) f(x,d\xi,\cdot) \m{ in } L^{1}_\eta((0,1)\times \R^{d(n+1)}).
\]
\end{lemma}

\begin{lemma}\label{ftau}
 For any $T\in \Tree$   there exists a transform $F\in \mathcal{F}$ of rank $n=|T|-1$ and the permutation $\sigma:\{1,\dots,n\}\rightarrow \{1,\dots,n\}$  such that    \[
 \tau(T,f)(t,x_1,\dots,x_n) = \izj F(f)(t,x_{\sigma(1)},\dots,x_{\sigma(n)},\eta)d\eta
 \]
 for any $f\in L^\infty((0,t_{*});L^{\infty}_\xi L^1_\eta ( (L^1\cap L^{\infty})(\R^d)))\cap L^\infty((0,t_{*});L^{\infty}_\eta L^1_\xi ( (L^1\cap L^{\infty})(\R^d)))$.
 \end{lemma}

\begin{lemma}\label{apro}
Let $f\in L^\infty((0,t_{*});\W((L^1\cap L^{\infty})(\R^d)))$ and $f^\ve$ denotes the approximation by convolution with periodic mollifying kernel in $\xi$- variable given in Lemma \ref{hmj}. For every $n-$rank transform  $F \in \mathcal{F}$ we have $ F(f^\ve)(t) \rightarrow  F(f)(t)$ in $L^{1}((0,1)\times\R^{dn})$ for almost all $t \in (0,t_{*})$ as $\ve \rightarrow 0$.
\end{lemma}

\begin{lemma}\label{taustar}
Let $f \in \W((L^{1}\cap L^{\infty})(\R^d))$, $T \in \Tree$ and $f^\ve$ be the approximation from Lemma \ref{hmj}. For every $m \in \{1,\dots,|T|\}$ we may represent $\tau(T,f^{\ve})$ as follows
\[
\tau(T,f^{\ve})(x_1,\dots,x_n) = \izj \tau_{*}^m(T,f^{\ve})(x_1,\dots,x_n,\xi_m)d\xi_m,
\]
\[
\tau_{*}^m(T,f^{\ve})(x_1,\dots,x_n,\xi_m):=\int_{[0,1]^{|T|-1}}\prod_{(k,l)\in E(T)}f^{\ve}(x_{l-1},\xi_k,\xi_l) \prod_{i\neq m} d\xi_i.
\]
Then as $\ve \rightarrow 0$
\[
\tau^m_{*}(T,f^{\ve}) \rightarrow \tau_{*}^m(T,f) \in L^{1}((0,1)\times \R^{dn}) \m{ and } \norm{\tau^m_{*}(T,f^\ve)}_{L^{\infty}((0,1)\times \R^{dn})} \leq \norm{f}_{\W(R^d)}^n. 
\]
\end{lemma}

\begin{proof}[Proof of Lemma \ref{taustar}]
We will proceed by induction on the number of vertices in a tree, but to give more insights in the first step we will discuss also three vertex trees. At first we take the only possible two vertex tree $T_2$. Here, for simplicity we denote $\xi_1 = \xi, \xi_2 = \eta, x_1 = x$. Then in case $m = 1$, $\tau^1_{*}$ actually coincides with $\tilde{\tau}$ from the proof of Lemma \ref{ftau} and the convergence follows from Proposition \ref{apro}. Anyway, let us repeat the short reasoning
\[
\tau_{*}^1(T_2,f^{\ve})(x,\xi) = F_{0}(f^{\ve})(x,\xi) = \izj  \rho^\ve(\xi-\xi')\izj f(x,\xi',d\eta) d\xi' = \rho^\ve * F_0(f)(x,\xi)
\]
and hence $\tau_{*}^1(T_2,f^{\ve}) \rightarrow F_0(f) = \tau^1_{*}(T_2,f)$ in $L^{1}((0,1)\times \R^d)$  with 
\[
 \norm{\tau^1_{*}(T_2,f^{\ve})}_{L^{\infty}((0,1)\times \R^d)} \leq \norm{\tau^1_{*}(T_2,f)}_{L^{\infty}((0,1)\times \R^d)} \leq \norm{f}_{\W(L^{\infty}(\R^d))}.
\]
In the case $m=2$ we have
\[
\tau_{*}^1(T_2,f^{\ve})(x,\eta) =\izj f^{\ve}(x,\xi,\eta)d\xi= \izj  \izj \rho^\ve(\xi-\xi') f(x,\xi',\eta) d\xi'd\xi = \izj f(x,\xi',\eta)d\xi'
\]
and since we have not only convergence but equality the claim follows in this case as well. 
Let us now discuss two three vertex tree: $T_3^1, T_3^2$. Let us begin with the tree with the root and two leaves  $T_3^1$. If $m=1$ (root), then again our $\tau_*$ coincides with $\tilde{\tau}$ from the proof of Lemma \ref{ftau}. Anyway, we may write
\begin{align*}
\tau_{*}^1(T_3^1,f^{\ve})(x_1,x_2,\xi_1) &= \izj f^\ve(x_1,\xi_1,\xi_2)d\xi_2 \izj f^{\ve}(x_2,\xi_1,\xi_3)d\xi_3
\\
&= \left(\rho^{\ve}* \izj f(x_1,\cdot,d\eta)\right)(\xi_1)\left(\rho^{\ve}* \izj f(x_2,\cdot,d\xi_3)\right)(\xi_1).
\end{align*}
Both terms on the RHS converge in $L^{1}((0,1)\times (\R^d))$ in $\xi_1$ and are uniformly bounded  by $\norm{f}_{\W(L^{\infty}(\R^d))}$, hence $\tau^1_{*}(T_3^1,f^{\ve}) \rightarrow \tau^1_{*}(T_3^1,f)$ in $L^{1}((0,1)\times \R^{2d})$ with the estimate
\[
 \norm{\tau^1_{*}(T^1_3,f^{\ve})}_{L^{\infty}((0,1)\times \R^{2d})} \leq \norm{\tau^1_{*}(f)}_{L^{\infty}((0,1)\times \R^{2d})} \leq \norm{f}^2_{\W(L^{\infty}(\R^d))}.
\]
If $m=2$, (a leaf) then we have
\begin{align*}
\tau_{*}^2(T_3^1,f^{\ve})(x_1,x_2,\xi_2) &= \izj f^\ve(x_1,\xi_1,\xi_2) \izj f^{\ve}(x_2,\xi_1,\xi_3)d\xi_3 d\xi_1\\
&= \izj f^\ve(x_1,\xi_1,\xi_2)\left(\rho^{\ve}* \izj f(x_2,\cdot,d\xi_3)\right)(\xi_1)d
\xi_1
\end{align*}
and we obtain the desired convergence using by Lemma \ref{hmj} with $g^{\ve}=\rho^{\ve}* \izj f(x_2,\cdot,d\xi_3)$. Furthermore, by Lemma \ref{impoest} and properties of convolution with mollifier 
\[
\norm{\tau^2_{*}(T^1_3,f^{\ve})}_{L^{\infty}((0,1)\times \R^{2d})} \leq \norm{f}^2_{\W(L^{\infty}(R^2d))}.
\]
If $m=3$ the argument is analogous.
Let us now discuss the tree $T^2_3$ with only one leaf  and three cases $m=1,2,3$. In case $m=1$ (a root) we have
\begin{align*}
\tau^1_{*}(T_3^2,f^{\ve})(x_1,x_2,\xi_1) &=  \izj f^\ve(x_1,\xi_1,\xi_2) \izj f^{\ve}(x_2,\xi_2,\xi_3)d\xi_3 d\xi_2
\\
&= \izj f^\ve(x_1,\xi_1,\xi_2)\left(\rho^{\ve}* \izj f(x_2,\cdot,d\xi_3)\right)(\xi_2)d
\xi_2
\end{align*}
and we obtain the convergence in this case again by Lemma \ref{hmj}.
In case $m=2$ (neither a root, nor a leaf) we have
\begin{align*}
\tau^2_{*}(T_3^2,f^{\ve})(x_1,x_2,\xi_2) &=  \izj f^\ve(x_1,\xi_1,\xi_2) d\xi_1 \izj f^{\ve}(x_2,\xi_2,\xi_3)d\xi_3 \\
&=\izj f(x_1,\xi_1,\xi_2) d\xi_1 \left(\rho^{\ve} *\izj f(x_2,\cdot,\xi_3)d\xi_3\right)(\xi_2).
\end{align*}
Since the convolution converges in $L^1$ and the integral $\izj f(x_1,\xi_1,\cdot) d\xi_1$ is bounded we obtain desired $L^1$ - convergence.
Finally, in case $m=3$ (a leaf) we get
\[
\tau^3_{*}(T_3^2,f^{\ve})(x_1,x_2,\xi_3) =  \izj\izj f^\ve(x_1,\xi_1,\xi_2) d\xi_1  f^{\ve}(x_2,\xi_2,\xi_3)d\xi_2 =\izj \izj f(x_1,\xi_1,\xi_2) d\xi_1 f^{\ve}(x_2,\xi_2,\xi_3)d\xi_2 
\]
and arrive at the desired $L^1$ - convergence applying Lemma \ref{hmj} to a constant sequence $g^\ve = \izj f(x_1,\xi_1,\cdot) d\xi_1$.
In all three cases $m=1,2,3$, we obtain the bound 
\[
\norm{\tau^m_{*}(T^2_3,f^{\ve})}_{L^{\infty}((0,1)\times R^{2d})} \leq \norm{f}^2_{\W(L^{\infty}(R^2d))}
\]
applying the properties of convolution with mollifier and in the case $m=3$ also Lemma \ref{impoest}.

Let us now assume that the claim holds for every $T \in \Tree_{\tilde{n}}$ for every $\tilde{n} \leq n$ and we take arbitrary $T \in \Tree_{n+1}$. We choose arbitrary $m \in \{1,\dots,|T|\}$. We denote by $T_m$ a subtree of $T$ rooted at vertex $m$.  If $m=1$ (a root) the claim follows from Lemma \ref{apro}. Otherwise we denote by $l_0$ the vertex that connects subtree $T_m$ with the rest of a tree, denoted by $T\setminus T_m$. Then
\[
\tau^m_{*}(T,f^{\ve})(x_1,\dots,x_{n},\xi_m) = \izj \tau^{l_0}(T \setminus T_m,f^\ve)(\xi_{l_0})f^{\ve}(x_{m-1},\xi_{l_0},\xi_m)dl_0 \tau_{*}^m(T_m,f^\ve)(\xi_m),
\]
where we omit some spatial indexes for clarity and we set $\tau_{*}^m(T_m,f^\ve)(\xi_m) \equiv 1$ in case when $m$ is a leaf. From the induction hypothesis $\tau_{*}^m(T_m,f^\ve) \rightarrow \tau_{*}^m(T_m,f)$ in $L^{1}$ and is uniformly bounded in $\ve$. Analogously $\tau^{l_0}(T \setminus T_m,f^\ve)\rightarrow \tau^{l_0}(T \setminus T_m,f)$ in $L^{1}$. Applying Lemma \ref{hmj} we arrive at $\tau^m_{*}(T,f^{\ve}) \rightarrow \tau^m_{*}(T,f)$ in $L^1((0,1)\times \R^{dn})$ and the estimate $\norm{\tau^m_{*}(T,f)}_{L^{\infty}((0,1)\times \R^{dn})} \leq \norm{f}^n_{\W(L^{\infty}(\R^d))}$ follows from the induction hypothesis and Lemma \ref{impoest}. Thus, we obtain the claim in case $T \in \Tree_{n+1}$ and we finish the proof applying mathematical induction principle. 
\end{proof}

\begin{thebibliography}{99}
\bibitem[JPS202]{JPS2025} Jabin P. E., Poyato D., Soler J., \textit{ Mean‐field limit of non‐exchangeable systems}, Communications on Pure and Applied Mathematics, 78(4), 651-741, (2025). %
\end{thebibliography}
\end{document}
